\documentclass[10pt,a4paper]{amsart}
\usepackage[margin=19mm]{geometry}
\usepackage{amsmath,amssymb,mathtools}
\usepackage[scr=rsfs,cal=euler]{mathalfa}
\usepackage{xfrac}
\usepackage[unicode,hidelinks,bookmarks=false]{hyperref}
\newcommand{\ie}{i.e.\ }
\newcommand{\cf}{cf.\ }
\newcommand{\resp}{respectively\ }
\newcommand{\Subsection}{\subsection}
\providecommand{\nobreakdash}{\nobreak\hbox{-}}
\setlength{\emergencystretch}{2em}
\multlinegap0pt
\usepackage{bm}
\usepackage{bbm}
\usepackage{dsfont}
\usepackage{xspace}
\usepackage{cancel}
\usepackage{mathtools}
\usepackage{amsthm}
\makeatletter
\@ifundefined{lemma}{\newtheorem{lemma}{Lemma}}{}
\makeatother
\title[On balanced Hermitian threefolds with parallel Bismut torsio]{On balanced Hermitian threefolds with parallel Bismut torsion}
\author{Quanting Zhao, Fangyang Zheng}
\date{Source: arXiv:2506.15141v2 (CC BY 4.0)}
\begin{document}
\maketitle

\noindent\emph{Source and licence.} On balanced Hermitian threefolds with parallel Bismut torsion. Version arXiv:2506.15141v2 (CC BY 4.0), distributed under the Creative Commons Attribution 4.0 International licence, as recorded in the arXiv licence metadata for this exact version and shown on the abstract page https://arxiv.org/abs/2506.15141v2. Full rights notice: licenses/arxiv-2506.15141-CC-BY-4.0.txt.\par\smallskip

\noindent\textit{Editorial scope.} Counted unit: the Lemma whose source label is \texttt{None} in the article (source0.tex lines 1086--1088, source label \texttt{None}), delivered together with the article's own complete original proof of it (source0.tex lines 1090--1226, one \texttt{proof} environment of 137 lines). No mathematics is written, repaired or re-derived: statement and proof are the article's own text line for line. Required context retained uncounted: none. Every definition, display and label of those lines is retained unchanged. Omitted from this package and not redistributed: 1--1085, 1089--1089, 1227--1951. Nothing inside a retained excerpt is omitted or altered: every retained body is delivered exactly as the article's lines are written, so no omission receipt of any kind accompanies this package. Citations: the retained excerpts carry no citation command at all, so no bibliography entry of the article is transcribed and no reference is redistributed. Reference adaptation: none was needed and none was made; every reference inside the delivered unit resolves inside it, so no unresolved reference or citation is produced. Printed slips kept literal: no printed word, symbol, hypothesis or number of the counted statement or of its proof was altered. Layout and class: the article's own class is \texttt{amsart} with packages including bbm, bbding, euscript, graphicx, amsthm, calc, mathrsfs, dsfont, CJK, fancyhdr, amscd, anysize, amsmath, amssymb; the wrapper is the lane's pinned \texttt{amsart} editorial wrapper at 10pt on A4, which restates the theorem-like environments the retained text needs and adds the article's own macro definitions that the retained text needs (none), with extra wrapper packages (bm, bbm, dsfont, xspace, cancel, mathtools). The delivered numbering is the unit's own. Assets: no retained span contains a figure environment, a graphics-inclusion command, a PDF-page inclusion, a table or a TikZ picture; no image file is copied into this package, the package declares no asset dependency and no image file is redistributed with it. Licence: version arXiv:2506.15141v2 of this article is distributed under the Creative Commons Attribution 4.0 International licence, as recorded in the arXiv licence metadata for this exact version and shown on the abstract page https://arxiv.org/abs/2506.15141v2. A full rights notice is delivered at licenses/arxiv-2506.15141-CC-BY-4.0.txt. Characters: every retained excerpt is pure ASCII, so the delivered engine, XeLaTeX with its default Latin Modern text font, reports no missing character.
\begin{lemma}
For any $s,s'\neq0$, the Lie algebras ${\mathfrak c}_s \cong {\mathfrak c}_{s'}$ if and only if $s=s'$ or $s'=\frac{1}{s}$.
\end{lemma}

\begin{proof}
Let $\{\varepsilon_i\}_{i=1}^6$ be the `standard' basis of the Lie algebra $\mathfrak{c}_s$ described in the above paragraphes. Consider a new basis $\{\tilde{\varepsilon}_i\}_{i=1}^6$ of $\mathfrak{c}_s$ given by
\[\tilde{\varepsilon}_1=\varepsilon_3,\quad \tilde{\varepsilon}_2=\varepsilon_4,\quad \tilde{\varepsilon}_3=\varepsilon_1,\quad \tilde{\varepsilon}_4=\varepsilon_2,\quad
\tilde{\varepsilon}_5=\frac{\varepsilon_5}{s},\quad \tilde{\varepsilon}_6=-\varepsilon_6.\]
It is not difficult to see that this new basis is a `standard' basis for the Lie algebra $\mathfrak{c}_{\frac{1}{s}}$. So ${\mathfrak c}_s$ and $\mathfrak{c}_{\frac{1}{s}}$ are isomorphic.
%So we have shown that
%when $s=s'$ or $s'=-s=\pm 1$ or $s'=\frac{1}{s}$ holds, the two Lie algebra ${\mathfrak c}_s$ and ${\mathfrak c}_{s'}$ are isomorphic.
Conversely, suppose that $f:{\mathfrak c}_s \rightarrow {\mathfrak c}_{s'}$ is an isomorphism with $s \neq s'$, which can be formulated as
\[f\begin{pmatrix} \varepsilon_1 \\ \varepsilon_2 \\ \varepsilon_3 \\ \varepsilon_4 \\ \varepsilon_6 \\ \varepsilon_5 \end{pmatrix}
=\begin{pmatrix} a_{11} & a_{12} & a_{13} & a_{14} & b_1 & 0\\
                 a_{21} & a_{22} & a_{23} & a_{24} & b_2 & 0\\
                 a_{31} & a_{32} & a_{33} & a_{34} & b_3 & 0\\
                 a_{41} & a_{42} & a_{43} & a_{44} & b_4 & 0\\
                    0    &     0   &     0   &     0   &  c  & 0 \\
                 d_{1}  & d_{2}  & d_{3}  & d_{4}  & d_6 & d\\
                 \end{pmatrix}
\begin{pmatrix} \varepsilon'_1 \\ \varepsilon'_2 \\ \varepsilon'_3 \\ \varepsilon'_4 \\ \varepsilon'_6 \\ \varepsilon'_5 \end{pmatrix},\]
where $\det(a_{ij})_{1\leq i,j \leq 4} \neq 0$, $c \neq 0$ and $d \neq 0$. The block pattern for the matrix above is due to the fact that  $f({\mathfrak c}_s') = {\mathfrak c}_{s'}'$ and $f({\mathfrak c}_s'') = {\mathfrak c}_{s'}''$. As $f$ preserves the Lie bracket, we have
$-[f(\varepsilon_1),f(\varepsilon_2)]=[f(\varepsilon_3),f(\varepsilon_4)]=f(\varepsilon_6)$, which implies
\begin{equation}\label{bck_12346}
c= \begin{vmatrix} a_{11} & a_{12} \\ a_{21} & a_{22} \\ \end{vmatrix} - \begin{vmatrix} a_{13} & a_{14} \\ a_{23} & a_{24} \end{vmatrix}
=\begin{vmatrix} a_{33} & a_{34} \\ a_{43} & a_{44} \\ \end{vmatrix} - \begin{vmatrix} a_{31} & a_{32} \\ a_{41} & a_{42} \end{vmatrix}.
\end{equation}
By $[f(\varepsilon_1),f(\varepsilon_5)]=-s f(\varepsilon_2)$ and $[f(\varepsilon_2),f(\varepsilon_5)]=sf(\varepsilon_1)$, we get the equalities (I):
\begin{align*}
a_{21}&=-\frac{a_{12}ds'}{s},& a_{22}&=\frac{a_{11}ds'}{s},& a_{23}&=-\frac{a_{14}d}{s}, & a_{24}&=\frac{a_{13}d}{s},
& b_2&=\frac{1}{s}\left(\begin{vmatrix}a_{11} & a_{12} \\ d_1 & d_2 \end{vmatrix}-\begin{vmatrix}a_{13} & a_{14} \\ d_3 & d_4 \end{vmatrix}\right),\\
a_{11}&=\frac{a_{22}ds'}{s},& a_{12}&=-\frac{a_{21}ds'}{s}, & a_{13}&=\frac{a_{24}d}{s},& a_{14}&=-\frac{a_{23}d}{s},
& b_1&=\frac{1}{s}\left(\begin{vmatrix}a_{23} & a_{24} \\ d_3 & d_4 \end{vmatrix}-\begin{vmatrix}a_{21} & a_{22} \\ d_1 & d_2 \end{vmatrix}\right).
\end{align*}
Similarly, by $[f(\varepsilon_3),f(\varepsilon_5)]=- f(\varepsilon_4)$ and $[f(\varepsilon_4),f(\varepsilon_5)]=f(\varepsilon_3)$ we get
the equalities (II):
\begin{align*}
a_{41}&=-a_{32}ds',& a_{42}&=a_{31}ds',& a_{43}&=-a_{34}d, & a_{44}&=a_{33}d,
& b_4&=\begin{vmatrix}a_{31} & a_{32} \\ d_1 & d_2 \end{vmatrix}-\begin{vmatrix}a_{33} & a_{34} \\ d_3 & d_4 \end{vmatrix},\\
a_{31}&=a_{42}ds',& a_{32}&=-a_{41}ds', & a_{33}&=a_{44}d,& a_{34}&=-a_{43}d,
& b_3&=\begin{vmatrix}a_{43} & a_{44} \\ d_3 & d_4 \end{vmatrix}-\begin{vmatrix}a_{41} & a_{42} \\ d_1 & d_2 \end{vmatrix}.
\end{align*}
Furthermore, by $[f(\varepsilon_1),f(\varepsilon_3)]=0$ and $[f(\varepsilon_2),f(\varepsilon_3)]=0$ we obtain
\begin{equation}\label{bck_1323}
\begin{vmatrix}a_{13} & a_{14} \\ a_{33} & a_{34} \end{vmatrix}=\begin{vmatrix}a_{11} & a_{12} \\ a_{31} & a_{32} \end{vmatrix},\quad
\begin{vmatrix}a_{23} & a_{24} \\ a_{33} & a_{34} \end{vmatrix}=\begin{vmatrix}a_{21} & a_{22} \\ a_{31} & a_{32} \end{vmatrix}.
\end{equation}
Similarly, $[f(\varepsilon_1),f(\varepsilon_4)]=0$ and $[f(\varepsilon_2),f(\varepsilon_4)]=0$ are equivalent to
\begin{equation}\label{bck_1424}
\begin{vmatrix}a_{13} & a_{14} \\ a_{43} & a_{44} \end{vmatrix}=\begin{vmatrix}a_{11} & a_{12} \\ a_{41} & a_{42} \end{vmatrix},\quad
\begin{vmatrix}a_{23} & a_{24} \\ a_{43} & a_{44} \end{vmatrix}=\begin{vmatrix}a_{21} & a_{22} \\ a_{41} & a_{42} \end{vmatrix}.
\end{equation}
Note that the first four columns of equalities (I) and (II) give us
\begin{eqnarray}
&& ( \frac{d^2s'^2}{s^2} -1)  \begin{pmatrix}a_{11} & a_{12} \\ a_{21} & a_{22} \end{pmatrix} \ \ =  \ \   (\frac{d^2}{s^2} -1 )  \begin{pmatrix}a_{13} & a_{14} \\ a_{23} & a_{24} \end{pmatrix} \ \ =  \ \  0,   \label{eq:5.19}\\
&& (d^2s'^2 -1)  \begin{pmatrix}a_{31} & a_{32} \\ a_{41} & a_{42} \end{pmatrix} \ \ =  \ \ ( d^2 -1)  \begin{pmatrix}a_{33} & a_{34} \\ a_{43} & a_{44} \end{pmatrix} \ \ =  \ \ 0 ,  \label{eq:5.20} \\
&& \begin{vmatrix} a_{11} & a_{12} \\ a_{21} & a_{22} \\ \end{vmatrix} \ \, = \ \, \frac{ds'}{s}(a_{11}^2 + a_{12}^2), \ \ \ \ \ \begin{vmatrix} a_{13} & a_{14} \\ a_{23} & a_{24} \\ \end{vmatrix} \ \, = \ \, \frac{d}{s}(a_{13}^2 + a_{14}^2), \label{eq:5.21} \\
&& \begin{vmatrix} a_{31} & a_{32} \\ a_{41} & a_{42} \\ \end{vmatrix} \ \, = \ \, ds' (a_{31}^2 + a_{32}^2), \ \ \ \ \ \begin{vmatrix} a_{33} & a_{34} \\ a_{43} & a_{44} \\ \end{vmatrix} \ \, = \ \, d(a_{33}^2 + a_{34}^2).  \label{eq:5.22}
\end{eqnarray}
We will divide the discussion into two cases depending on whether
the matrix $\begin{pmatrix}a_{33} & a_{34} \\ a_{43} & a_{44} \end{pmatrix} $ is zero or not.

\noindent {\bf Case 1.} If $\begin{pmatrix}a_{33} & a_{34} \\ a_{43} & a_{44} \end{pmatrix} \neq 0$.

In this case the second equality of (\ref{eq:5.20}) implies that $d^2=1$
%and $\begin{vmatrix}a_{33} & a_{34} \\ a_{43} & a_{44} \end{vmatrix}=d(a_{33}^2 + a_{34}^2) \neq 0$.
Let us assume $d=1$ as the case $d=-1$ can be argued similarly. Hence $\begin{vmatrix}a_{33} & a_{34} \\ a_{43} & a_{44} \end{vmatrix} = d(a_{33}^2+a_{34}^2)>0$. We make the following claim:
\vspace{0.2cm}

\noindent {\em Claim 1:}  $\begin{pmatrix}a_{13} & a_{14} \\ a_{23} & a_{24} \end{pmatrix} \neq 0$
and $\begin{pmatrix}a_{31} & a_{32} \\ a_{41} & a_{42} \end{pmatrix} \neq 0$.

If on the contrary that $\begin{pmatrix}a_{13} & a_{14} \\ a_{23} & a_{24} \end{pmatrix}=0$. Then we have $
\begin{vmatrix}a_{11} & a_{12} \\ a_{21} & a_{22} \end{vmatrix} \neq 0$ since $\det(a_{ij}) \neq 0$.
By the first equality of (\ref{eq:5.19}) we conclude that $(\frac{ds'}{s})^2=(\frac{s'}{s})^2=1$, hence $s'=-s$ as we have assumed that  $s \neq s'$ at the very beginning. The first equality of (\ref{eq:5.21}) gives us
$$
\begin{vmatrix}a_{11} & a_{12} \\ a_{21} & a_{22} \end{vmatrix}=  -(a_{11}^2+a_{12}^2) <0.
$$
On the other hand, because $a_{13}=a_{14}=a_{23}=a_{24}=0$,  the equalities \eqref{bck_1323} and \eqref{bck_1424} become
\[\begin{vmatrix}a_{11} & a_{12} \\ a_{31} & a_{32} \end{vmatrix}=\begin{vmatrix}a_{21} & a_{22} \\ a_{31} & a_{32} \end{vmatrix}=0,\quad \quad \quad
\begin{vmatrix}a_{11} & a_{12} \\ a_{41} & a_{42} \end{vmatrix}=\begin{vmatrix}a_{21} & a_{22} \\ a_{41} & a_{42} \end{vmatrix}=0.\]
Since vectors $(a_{11},a_{12})$ and $(a_{21},a_{22})$ are linearly independent, we conclude that
$\begin{pmatrix}a_{31} & a_{32} \\ a_{41} & a_{42} \end{pmatrix} =0$. But then
the equality \eqref{bck_12346} would give us
\[c=\begin{vmatrix}a_{11} & a_{12} \\ a_{21} & a_{22} \end{vmatrix}=-(a_{11}^2+a_{12}^2)<0,\quad \quad
c=\begin{vmatrix} a_{33} & a_{34} \\ a_{43} & a_{44} \end{vmatrix}=a_{33}^2+a_{34}^2>0,\]
which is a contradiction. So we must have $\begin{pmatrix}a_{13} & a_{14} \\ a_{23} & a_{24} \end{pmatrix} \neq 0$. Similarly, we must have
 $\begin{pmatrix}a_{31} & a_{32} \\ a_{41} & a_{42} \end{pmatrix} \neq 0$ as well, and Claim 1 is proved.
\vspace{0.2cm}

\noindent {\em Claim 2:}  $\begin{pmatrix}a_{11} & a_{12} \\ a_{21} & a_{22} \end{pmatrix} \neq 0$.

Assume on the contrary that the above matrix is zero. Then by \eqref{bck_1323} and \eqref{bck_1424} we get
$$ \begin{vmatrix} a_{13} & a_{14} \\ a_{33} & a_{34} \\ \end{vmatrix} = \begin{vmatrix} a_{13} & a_{14} \\ a_{43} & a_{44} \\ \end{vmatrix} = \begin{vmatrix} a_{23} & a_{24} \\ a_{33} & a_{34} \\ \end{vmatrix} = \begin{vmatrix} a_{23} & a_{24} \\ a_{43} & a_{44} \\ \end{vmatrix} =0.
$$
Since the vectors $(a_{33},a_{34})$ and $(a_{43},a_{44})$ are linearly independent, we conclude that $\begin{pmatrix}a_{13} & a_{14} \\ a_{23} & a_{24} \end{pmatrix} =0$, contradicting to Claim 1. So we must have  $\begin{pmatrix}a_{11} & a_{12} \\ a_{21} & a_{22} \end{pmatrix} \neq 0$, and Claim 2 is proved.

By the claims and the equalities (\ref{eq:5.19}), (\ref{eq:5.20}), we conclude that $s^2=s'^2=1$. Since $s\neq s'$, we know that $(s,s')$ is either $(1,-1)$ or $(-1,1)$. The equality \eqref{bck_12346} now takes the following form
\[\begin{split}
c &=\begin{vmatrix} a_{11} & a_{12} \\ a_{21} & a_{22} \\ \end{vmatrix} - \begin{vmatrix} a_{13} & a_{14} \\ a_{23} & a_{24} \end{vmatrix}
=-(a^2_{11}+a_{12}^2)-\frac{1}{s}(a^2_{13}+a^2_{14}), \\
c &= \begin{vmatrix} a_{33} & a_{34} \\ a_{43} & a_{44} \\ \end{vmatrix} - \begin{vmatrix} a_{31} & a_{32} \\ a_{41} & a_{42} \end{vmatrix}
=(a_{33}^2 + a_{34}^2) -s'(a_{31}^2+a_{32}^2).
\end{split}\]
If $(s,s')=(1,-1)$, then the above two formula would dictate that $c<0$ and $c>0$ at the same time, which is absurd. So we must have $s=-1$ and $s'=1$.
Next let us write
\[r_1 = \sqrt{a_{11}^2 + a_{12}^2},\quad r_2 = \sqrt{a_{13}^2+a_{14}^2},\quad r_3=\sqrt{a_{31}^2+a_{32}^2},\quad r_4=\sqrt{a_{33}^2+a_{34}^2},\]
where $r_1,r_2,r_3,r_4$ are positive numbers. Then the four blocks of the matrix $(a_{ij})_{1\leq i,j\leq 4}$ can be reformulated as
\begin{align*}
\begin{pmatrix} a_{11} & a_{12} \\ a_{21} & a_{22}\end{pmatrix} &= \begin{pmatrix} a_{11} & a_{12} \\ a_{12} & -a_{11} \end{pmatrix}
=r_1 \begin{pmatrix} p_1 & p_2 \\ p_2 & -p_1 \end{pmatrix}, &
\begin{pmatrix} a_{13} & a_{14} \\ a_{23} & a_{24} \end{pmatrix} &= \begin{pmatrix} a_{13} & a_{14} \\ a_{14} & -a_{13} \end{pmatrix}
=r_2 \begin{pmatrix} p_3 & p_4 \\ p_4 & -p_3 \end{pmatrix}, \\
\begin{pmatrix} a_{31} & a_{32} \\ a_{41} & a_{42} \end{pmatrix} &= \begin{pmatrix} a_{31} & a_{32} \\ -a_{32} & a_{31} \end{pmatrix}
=r_3 \begin{pmatrix} p_5 & p_6 \\ -p_6 & p_5 \end{pmatrix}, &
\begin{pmatrix} a_{33} & a_{34} \\ a_{43} & a_{44} \end{pmatrix} &= \begin{pmatrix} a_{33} & a_{34} \\ -a_{34} & a_{33} \end{pmatrix}
=r_4 \begin{pmatrix} p_7 & p_8 \\ -p_8 & p_7 \end{pmatrix},
\end{align*}
where $p_1^2+p_2^2=p_3^2+p_4^2=p_5^2+p_6^2=p_7^2+p_8^2 =1$. The equality \eqref{bck_1323} then implies that
\[r_1 r_3(p_1 p_6 - p_2 p_5)=r_2 r_4(p_3 p_8 - p_4 p_7 ),\quad r_1 r_3(p_2 p_6 + p_1 p_5)=r_2 r_4(p_4 p_8 + p_3 p_7 ). \]
Square both equalities above and add them up,  we get $r^2_1 r^2_3 = r^2_2 r^2_4$,
which leads to $r_1 r_3 = r_2 r_4$. On the other hand, the equality \eqref{bck_12346} now says that
\[ c=r_2^2 -r_1^2 = r_4^2 -r_3^2.\]
Note that the Lie algebra isomorphism $f$ maps the center to the center and $c$ is the coefficient there, so $c\neq 0$. If $c>0$, then $r_2>r_1>0$ and $r_4>a_3>0$, which would lead to $r_2r_4>r_1r_3$. Similarly when $c<0$ we get $r_2r_4<r_1r_3$. This contradiction shows that our Case 1 cannot occur.

\vspace{0.2cm}

\noindent {\bf Case 2.} If $\begin{pmatrix}a_{33} & a_{34} \\ a_{43} & a_{44} \end{pmatrix} = 0$.

In this case by $\det(a_{ij})\neq 0$ we get $\begin{vmatrix}a_{13} & a_{14} \\ a_{23} & a_{24} \end{vmatrix} \neq 0$ and $
\begin{vmatrix}a_{31} & a_{32} \\ a_{41} & a_{42} \end{vmatrix} \neq 0$.
The equalities (\ref{eq:5.19}) and (\ref{eq:5.20}) imply that  $(ds')^2=1$ and $(\frac{d}{s})^2=1$, that is, $d^2=s^2=\frac{1}{s'^2}$.
The equalities \eqref{bck_1323} and \eqref{bck_1424} now take the following form
\[\begin{vmatrix}a_{11} & a_{12} \\ a_{31} & a_{32} \end{vmatrix}=\begin{vmatrix}a_{21} & a_{22} \\ a_{31} & a_{32} \end{vmatrix}=0,\quad \quad
\begin{vmatrix}a_{11} & a_{12} \\ a_{41} & a_{42} \end{vmatrix}=\begin{vmatrix}a_{21} & a_{22} \\ a_{41} & a_{42} \end{vmatrix}=0.\]
Since the vectors $(a_{31}, a_{32})$ and $(a_{41},a_{42})$ are linearly independent to each other, the above equalities imply that
$\begin{pmatrix}a_{11} & a_{12} \\ a_{21} & a_{22} \end{pmatrix} =0$. Hence by the equality \eqref{bck_12346} we get
\[c=-\begin{vmatrix}a_{13} & a_{14} \\ a_{23} & a_{24} \end{vmatrix} = - \frac{d}{s}(a_{13}^2+a_{14}^2),
\quad c=-\begin{vmatrix}a_{31} & a_{32} \\ a_{41} & a_{42} \end{vmatrix} = -ds'(a_{31}^2+a_{32}^2).\]
This indicates that $\frac{d}{s}$ and $ds'$ have the same sign, which leads to $s'=\frac{1}{s}$. This completes the proof of the lemma.
\end{proof}

\end{document}
